\begin{figure}[t]
\centering
\begin{tikzpicture}[
  x=1cm,y=1cm,
  vertex/.style={circle,draw,fill=white,minimum size=5mm,inner sep=0pt},
  edge/.style={-{Stealth[length=2mm]},semithick},
  every node/.style={font=\small}
]
  \fill[black!4,rounded corners=5pt] (0.1,-1.65) rectangle (3.55,1.8);
  \node at (-2.8,1.45) {nonincreasing dimension};
  \node at (1.83,2.10) {equal dimension};

  \node[vertex] (S0) at (-4.4,0) {};
  \node[vertex] (S1) at (-2.8,0) {};
  \node[vertex] (S2) at (-1.2,0) {};
  \node[vertex] (A) at (0.7,0) {};
  \node[vertex] (B) at (2.7,1) {};
  \node[vertex] (C) at (2.7,-1) {};

  \node[above=5pt] at (S0.north) {$S_0$};
  \node[above=5pt] at (S1.north) {$S_1$};
  \node[above=5pt] at (S2.north) {$S_2$};
  \node[above left=5pt] at (A.north west) {$A$};
  \node[above=5pt] at (B.north) {$B$};
  \node[right=5pt] at (C.east) {$C$};

  \draw[edge] (S0) -- (S1);
  \draw[edge] (S1) -- (S2);
  \draw[edge] (S2) -- (A);
  \draw[edge] (A) -- (B);
  \draw[edge] (B) -- (C);
  \draw[edge] (C) -- (A);

  \node[align=center] at (-2.8,-0.95)
    {$S\longrightarrow S'$ implies $d(S')\leq d(S)$};
  \node[align=center] at (1.83,-2.10)
    {$A,B,C\in\mathcal C_{\mathrm{ex}}$\\[2pt]
     unique outgoing neighbor at each cycle vertex};
\end{tikzpicture}
\caption{A schematic path in the graph of simple objects, with an
illustrative three-vertex cycle. Dimension is nonincreasing along every
arrow; arrows before the cycle may also preserve dimension. On a cycle,
Lemma~\ref{lem:equality-descent} gives equality, membership in
$\mathcal C_{\mathrm{ex}}$, and a unique outgoing neighbor at every
vertex, so the path cannot leave the cycle.}
\label{fig:dimension-descent}
\end{figure}
